\section{AI 的参差性：天才 PhD 与10岁小孩}

\subsection{在轨道上是超级智能，脱轨是10岁小孩}

Karpathy 用了一个极其精准的比喻来描述当前 AI 的状态：

\begin{quote}
``I simultaneously feel like I'm talking to an extremely brilliant PhD student who's been a systems programmer for their entire life and a 10-year-old. This is so weird because humans wouldn't have that combination.''

——我同时感觉在跟一个极其聪明的终身系统程序员 PhD 和一个10岁小孩对话。这太奇怪了，因为人类不会出现这种组合。
\end{quote}

人类的能力更加"耦合"——各方面水平差不多。但 Agent 有远超人类的\textbf{参差性（jaggedness）}。有时候你让它实现一个功能，它回来的东西完全离谱，然后你们陷入错误循环，让人抓狂。

\subsection{为什么参差性存在：强化学习的盲区}

Karpathy 分析了根本原因：模型是通过强化学习（RL）训练的，所以它们能改进的只有\textbf{可验证的东西}——程序是否正确？单元测试是否通过？但更"软"的能力，比如理解你的意图、知道什么时候该追问，不在 RL 的优化范围内。

他用了一个直观的例子：你去问最先进的 ChatGPT 讲个笑话，知道你会得到什么笑话吗？

\begin{quote}
``Why do scientists not trust atoms? Because they make everything up. — This is the joke you would get three or four years ago and this is the joke you still get today.''
\end{quote}

模型在 Agent 任务上已经能连续跑几个小时、移山倒海，但讲笑话还是和五年前一模一样。因为笑话在强化学习的优化范围之外。

\begin{warningbox}{流行叙事的问题}
有一个流行假说：在可验证领域（编程、数学）变强就会在所有领域变强。Karpathy 认为这并没有发生，或者只发生了一点点，远不到令人满意的程度。"我们不会免费地在所有领域获得智能和能力"——有些领域在被优化，有些不在，而这一切都压缩在不透明的神经网络里。
\end{warningbox}

\subsection{一个模型还是一千个大脑：物种分化的缺失}

Sarah 问了一个她自称"有点亵渎"的问题：既然参差性持续存在，是否应该把单体模型拆开——拆成可以在不同领域分别优化的专业化版本？

Karpathy 认为当前的实验室在追求单一的\textbf{monoculture}——一个模型在所有领域都要聪明。但他觉得应该有更多\textbf{speciation（物种分化）}。动物界有极其多样化的大脑，有些动物的视觉皮层过度发达。AI 也应该如此：保留认知核心但特化到具体任务。

但他坦承目前还没有看到太多物种分化，原因有二：
\begin{enumerate}
    \item 实验室不知道用户会问什么，必须覆盖所有可能
    \item \textbf{操纵大脑的科学还不够成熟}——通过 context window 做定制简单又便宜，但真正触碰 weights（微调而不丢失能力、持续学习、领域特化）风险大得多
\end{enumerate}

\begin{importantbox}{Context Window vs Weights}
当前获得个性化的主要方式是通过 context window——简单、便宜、安全。但真正的物种分化需要触碰 weights，而动了 weights 就是在改变整个模型和它的智能。微调而不丢失能力、持续学习、领域特化——这些技术还在发展中，这是物种分化的技术瓶颈。
\end{importantbox}

\subsection{本章小结}

AI 的参差性（jaggedness）是当前阶段的核心特征：在可验证的 RL 优化领域表现出超级智能，在优化范围之外则停滞不前。这挑战了"通用智能免费涌现"的叙事。物种分化可能是解决方案，但受制于"操纵大脑的科学"尚未成熟。


